% LSPV-4 REPORT::555::2026-11-11 18:38 PST
\documentclass[11pt]{article}
\usepackage{amsmath,amssymb,booktabs,hyperref}
\begin{document}
\begin{center}
{\Large LSPV-4/555: Source-order retention versus reference-byte cost}\\
November 11, 2026, 18:38 PST\\
Time-Space Resource Auditor, historical Route A
\end{center}
\noindent\textbf{PROVED/RESTRICTED plus COMPUTATIONAL. P versus NP OPEN.}
\section*{Candidate and falsifiers}
Test whether moving terminal-dead but mandatory original SOURCE records before all paid Resolution deductions reduces retention cost at unchanged logical work. Falsifiers are native-checker rejection of the reordered proof, an unrecognized parent dependence, an abstract retention cost exceeding a concrete accepted cost, and invalid references crossing a decimal-digit boundary. The unchanged archived /386 SOURCE/paid-AND2/DEF/RES checker is Git blob 338305c0ce101704205e57543756b8127074a986. The prior W29 seed is blob 035d681e2d18e8b12191b984d59bcf870d3915d0.
\section*{Theorem, quantifiers and proof}
Let $t$ denote the number of disjoint original signed $K_4$ components with guard $g=6t+1$, odd charge at only the first vertex and exactly $16t+1$ mandatory original SOURCE records. The fixed proof has 28 paid RES and one unused but globally charged AND2 owner, so native work is $W=29$. In the LATE trace, 17 SOURCEs precede 27 RES steps, then $16(t-1)$ unrelated SOURCEs precede the final RES. In the FRONT trace all SOURCEs precede the same 28 RES steps, with all parent indices renumbered. The SOURCE specification, signed pivots, resolvents, and terminal empty clause are identical.

Under read-only SOURCEs, fresh paid output registers and last-use deletion, let $A$ count live paid RES result slot-steps sampled after each record. Then for these exact two trace orders,
\[
\boxed{A_{\rm late}=144+16(t-1),\qquad A_{\rm front}=144}.
\]
The original 27 nonterminal RES results each have exactly one later paid consumer. The late trace has 56 parent applications and 92 parent-reference decimal digits, comprising 29 SOURCE references (38 digits) and 27 paid references (54 digits). The front trace has exactly
\[
\boxed{D_{\rm front}=38+\sum_{k=1}^{27}\operatorname{dig}_{10}(16t+k)}
\]
reference digits. Only these parent-index digits change in compact JSON serialization; hence total certificate-byte difference equals $D_{\rm front}-92$. Both traces retain maximum fresh result slots 9 and dependency span 7. Every extra late SOURCE record adds one live sample of the previously derived $[-g]$, but moving it to the front adds none. At $t=3,4$, FRONT strictly dominates LATE in area with equal W, peak, references and certificate bytes. For increasing $t$, the area savings are linear while added reference digits grow logarithmically. This is not an optimum over all possible proof DAGs.
\section*{Native verification and hostile controls}
Exact published program experiments/arc/lspv4_555_source_prefix_area.cjs (blob 342d76f278265e2c008a2a5d6100703a621b5cae) recreates the prior seed byte-for-byte and independently verifies signed SOURCE and RES records, count-once owners, reference applications and live sets for t=1,2,3,4,5,8,12,16,32,64. The unchanged original archived checker body was also evaluated in hosted V8 on both traces for t=1,2,3,4,5,8,12,16 with virtual I/O only. Both variants ACCEPT for t=1,3,4,5,8,12,16 and each passes the original eleven bundled hostile tests; both t=2 variants REJECT with the known disconnected parity self-test error. Five additional forged SOURCE/guard, parent, cost and owner packets REJECT. Independent original Node OS subprocess execution is not claimed; the program includes that optional replay. Fake hosted digest placeholders are not cryptographic SHA values.

\begin{center}
\begin{tabular}{rrrrrr}
\toprule
$t$ & SOURCE & $A_{\rm late}$ & $A_{\rm front}$ & $D_{\rm front}$ & bytes LATE/FRONT\\
\midrule
3 &49&176&144&92&4994/4994\\
4 &65&192&144&92&6066/6066\\
8 &129&256&144&119&10356/10383\\
16&257&384&144&119&18932/18959\\
64&1025&1152&144&146&74818/74872\\
\bottomrule
\end{tabular}
\end{center}
Only signed structural replay, not native checker acceptance, was run at t=64. Published accepted front-packet blobs: t=3 8131078e841a1d0b69641534020134b902b75e75 and t=16 9e180273b63e5fc4190c43c58d03e3f2edc2a38c. Raw audit blob 0efe9bdc3864e59f6361e319633ad2698728fab0.

Every one of 16t vertex SOURCE clauses contains numeric literal -g. Changing any decimal digit of one of these literals while holding its SOURCE specification fixed invalidates mandatory exact-source checking. Therefore ordinary deterministic byte-probe exact-source verification on an accepting trace requires inspecting at least $16t\operatorname{dig}_{10}(6t+1)=\Omega(t\log t)$ sensitive digit positions. This concerns evidence verification, not decision-time hardness for arbitrary SAT algorithms. Physical checker RAM, parser buffers, computation discovery, retrieval and scratch remain separate costs. A schedule-insensitive abstract rule charging the late area for FRONT overestimates a genuine concrete cost by $16(t-1)$ and violates the required abstraction cost inequality; exact live-set transformers do not, for these restricted traces.

\section*{Outstanding obligations}
The W29 paid proof is a tree, not a genuinely multiply consumed paid shared DAG; a single unused owner is charged once. GF(2), paid extension, learning, branches, recomputation, compact affine rows and cross-rule context remain legal and unexcluded. Complete historical R_mix341 byte grammar and weights remain unauthenticated. Both root trackers (proof1 blob f46f6aae25571340a5fd17a5181d51c271509947; proof2 blob 42943fe447c6061fb8009a526f46ac8a6684765c), issue 3, AGENTS, CODEX, project skill, original checker, LSPV-4/549 and /552--554, 539--549 lineage and recent LSPV-1/585, LSPV-2/569, LSPV-3/582 and LSPV-5/498 were reviewed.
\par\noindent\textbf{TRACKERS REVIEWED; NO MATERIAL CHANGE.}
P2-D1 to P2-C3 to P2-B1 remain OPEN; P2-B2/B3 separately OPEN. P versus NP remains OPEN.

\begin{thebibliography}{9}
\bibitem{NH} J. Nordstrom and J. Hastad, Towards an Optimal Separation of Space and Length in Resolution, Theory of Computing 9 (2013), 471--557.
\bibitem{AD} A. Atserias and V. Dalmau, A Combinatorial Characterization of Resolution Width, JCSS 74 (2008), 323--334.
\bibitem{CR} S. Cook and R. Reckhow, The Relative Efficiency of Propositional Proof Systems, JSL 44 (1979), 36--50.
\end{thebibliography}
\end{document}
% END_OF_VERIFIED_REPORT::LSPV-4::555
